\chapter{Futures Contracts and Their Exchanges}\label{ch:m1:futures-contracts-and-exchanges}

With the index at 6\,000, one E-mini S\&P 500 future is an exposure to
\$300\,000 of American shares. It costs nothing to enter: nobody pays
\$300\,000, and the deposit the clearing house asks for is a small fraction of
it. The contract can be sold short as easily as bought, trades from Sunday
evening to Friday afternoon, and its smallest price step, a quarter of a
point, is worth \$12.50. Every one of those facts is a line in a rulebook,
and different for each of the thousands of futures listed around the world.
A firm that trades futures begins by turning rulebooks into a table. This
chapter builds the table; the next three explain how orders match, how the
deposit is computed, and how a futures price relates to the thing it is a
future on.

\section{The contract}

\begin{definition}[Futures contract]\label{def:m1:futures-contracts-and-exchanges:future}
A \emph{futures contract}\index{futures contract} is a standardised
agreement, listed by an exchange and guaranteed by its clearing house
(\cref{ch:m1:clearing-and-settlement}), to buy or sell a specified quantity of
an underlying at a specified future date, at a price agreed today. Gains and
losses are settled in cash every day; at expiry the contract is either
settled in cash against a reference price or by delivery of the underlying.
\end{definition}

Standardisation is the product. Two parties who want a forward on crude oil
must agree on grade, place, date, size, credit and documentation. The
exchange fixes all of these once, so that the only thing left to negotiate is
price, and every contract of a given month is interchangeable with every
other: a position bought from one counterparty can be closed by selling to
another.

\begin{definition}[Multiplier, tick and tick value]\label{def:m1:futures-contracts-and-exchanges:tick}
The \emph{contract multiplier}\index{contract multiplier} converts the quoted
price into money: the contract's notional value is price times multiplier.
The \emph{tick} is the minimum price increment, and the \emph{tick
value}\index{tick value} is tick times multiplier, the profit or loss of one
contract for one tick.
\end{definition}

\begin{definition}[Expiry month code and front month]\label{def:m1:futures-contracts-and-exchanges:code}
A futures symbol is a product root followed by an \emph{expiry month
code}\index{expiry month code}, one letter for the month (F, G, H, J, K, M, N,
Q, U, V, X, Z for January to December), and the last digit of the year. The
\emph{front month}\index{front month} is the nearest listed expiry, usually
the most traded.
\end{definition}

\begin{omfigure}
\begin{tikzpicture}[font=\small]
\node[font=\Huge\ttfamily,text=omDef] (r) at (0,0.2) {ES};
\node[font=\Huge\ttfamily,text=omThm] (m) at (1.7,0.2) {Z};
\node[font=\Huge\ttfamily,text=omExo] (y) at (3.1,0.2) {6};
\node[below=4pt of r,font=\footnotesize,align=center,text=omDef] {root:\\the product};
\node[below=4pt of m,font=\footnotesize,align=center,text=omThm] {month:\\December};
\node[below=4pt of y,font=\footnotesize,align=center,text=omExo] {year:\\2026};
\foreach \c/\n [count=\i] in {F/Jan,G/Feb,H/Mar,J/Apr,K/May,M/Jun,N/Jul,Q/Aug,U/Sep,V/Oct,X/Nov,Z/Dec}{
  \pgfmathsetmacro{\x}{4.2+0.78*\i}
  \pgfmathtruncatemacro{\q}{mod(\i,3)==0}
  \ifnum\q=1
    \node[draw=omThm,fill=omThm!12,minimum width=0.7cm,minimum height=0.7cm,font=\ttfamily\bfseries] at (\x,0.25) {\c};
  \else
    \node[draw=omExo,minimum width=0.7cm,minimum height=0.7cm,font=\ttfamily] at (\x,0.25) {\c};
  \fi
  \node[font=\scriptsize] at (\x,-0.4) {\n};}
\end{tikzpicture}

\omcaption{Reading a futures symbol. The shaded codes, H, M, U and Z, are the
quarterly cycle on which equity-index and bond futures are listed; energy
contracts list every month.}\label{fig:m1:futures-contracts-and-exchanges:symbol}
\end{omfigure}

\section{The exchange groups}

Futures liquidity is concentrated by product, and each product lives on one
exchange: unlike shares (\cref{ch:m1:us-equity-market-structure}), a future is
the intellectual property of the exchange that lists it and clears only at
that exchange's clearing house. A position opened on one exchange cannot be
closed on another. The consequence is a small number of groups, each a
near-monopoly in its products: in the United States one group owns the
Chicago and New York exchanges that list the benchmark equity-index,
interest-rate, energy, metal and agricultural contracts; an Atlantic group
lists Brent crude and European short-term rates in London and soft
commodities in New York; the German exchange group lists the euro area's
equity-index and government-bond futures; and national exchanges in Asia and
Latin America list their own indices, rates and commodities. Competition
between exchanges happens at the launch of a product, rarely afterwards.

\section{Six specifications}

\begin{center}
\small
\begin{tabular}{llllll}
\toprule
Contract & Unit & \omterm{def:m1:futures-contracts-and-exchanges:tick}{Tick} & \omterm{def:m1:futures-contracts-and-exchanges:tick}{Tick value} & Expiries & Settlement \\
\midrule
E-mini S\&P 500 (ES) & \$50 $\times$ index & 0.25 & \$12.50 & quarterly & cash \\
EURO STOXX 50 (FESX) & \euro10 $\times$ index & 1 & \euro10 & quarterly & cash \\
10-year T-note (ZN) & \$100\,000 face & $1/64$ & \$15.625 & quarterly & delivery \\
Euro-Bund (FGBL) & \euro100\,000 face & 0.01 & \euro10 & quarterly & delivery \\
WTI crude (CL) & 1\,000 barrels & \$0.01 & \$10 & monthly & delivery \\
Brent crude (B) & 1\,000 barrels & \$0.01 & \$10 & monthly & cash or EFP \\
\bottomrule
\end{tabular}
\end{center}

\begin{dated}{2026-09}{Details from the rulebooks}\label{dat:m1:futures-contracts-and-exchanges:specs}
\textbf{ES}: final settlement is in cash against a special opening quotation
of the index, computed from the opening prices of its component stocks on
the third Friday of the contract month; calendar spreads trade in increments
of 0.05. \textbf{FESX}: the twelve nearest quarterly months are listed;
trading ends at 12:00 CET on the third Friday and the final settlement price
is the average of the index between 11:50 and 12:00; regular trading runs
from 02:10 to 22:00 CET. \textbf{ZN}: deliverable notes have a remaining
maturity of at least six and a half and less than eight years; prices are
quoted in points and thirty-seconds. \textbf{FGBL}: a notional 6\% bond;
deliverable bonds have 8.5 to 10.5 years remaining. \textbf{CL}: delivery at
Cushing, Oklahoma. \textbf{B}: up to 156 consecutive months listed; trading
ends on the last business day of the second month before the contract month;
settlement by exchange for physical with an option to settle in cash against
an index.
\end{dated}

Bond futures are quoted in a base that predates computers. A price of
\texttt{112-165} means $112 + 16.5/32$: the first two digits after the dash
count thirty-seconds of a point and a third digit, 2, 5 or 7, a quarter, a
half or three quarters of a thirty-second. With a \omterm{def:m1:futures-contracts-and-exchanges:tick}{tick} of half a
thirty-second, $1/64$ of a point on \$100\,000 of face value, the \omterm{def:m1:futures-contracts-and-exchanges:tick}{tick value}
is \$15.625. Every system that touches these prices should hold them as
integers of \omterm{def:m1:futures-contracts-and-exchanges:tick}{ticks}: a price is then never ``almost'' on the grid, equality
tests are exact, and no P\&L depends on rounding. (One sixty-fourth happens
to be exact in binary floating point; 0.01, the \omterm{def:m1:futures-contracts-and-exchanges:tick}{tick} of three other contracts
in the table, is not.)

\begin{omfigure}
\begin{tikzpicture}
\begin{axis}[width=10cm,height=4.6cm,ybar,bar width=14pt,
  ylabel={one tick (bp of value)},
  symbolic x coords={ES,FESX,ZN,FGBL,CL,B},xtick=data,
  tick label style={font=\small},label style={font=\small},xticklabel style={font=\small\ttfamily\strut},
  ymin=0,ymax=2.3,ymajorgrids,grid style={gray!25},
  nodes near coords,nodes near coords style={font=\footnotesize,/pgf/number format/.cd,fixed,fixed zerofill,precision=2},
  table/col sep=comma]
\addplot[fill=omDef!60,draw=omDef] table[x=root,y=tick_bp]{figdata/markets-1/18-futures-contracts-and-exchanges/specs.csv};
\end{axis}
\end{tikzpicture}

\omcaption{The \omterm{def:m1:futures-contracts-and-exchanges:tick}{tick} as a fraction of the contract's value, at illustrative
prices (ES 6\,000, FESX 5\,400, ZN 112, FGBL 128, CL 70, B 74). A coarse \omterm{def:m1:futures-contracts-and-exchanges:tick}{tick}
means a wide minimum spread, long queues and a premium on being first; a fine
\omterm{def:m1:futures-contracts-and-exchanges:tick}{tick} means price competition. The matching consequences are the subject of
\cref{ch:m1:matching-algorithms-and-implied-spreads}. Data: the tutorial's
table.}\label{fig:m1:futures-contracts-and-exchanges:tickbp}
\end{omfigure}

\section{Sessions, settlement and limits}

\begin{definition}[Daily settlement price]\label{def:m1:futures-contracts-and-exchanges:settle}
The \emph{daily settlement price}\index{daily settlement price} is the price
published by the exchange each day at which every open position is marked
and \omterm{def:m1:clearing-and-settlement:margin}{variation margin} is computed. It is determined by a rule, typically a
volume-weighted average of trades in a short window, not by the last trade.
\end{definition}

\begin{definition}[Open interest]\label{def:m1:futures-contracts-and-exchanges:oi}
\emph{Open interest}\index{open interest} is the total of all contracts
entered into and not yet offset by a transaction, by delivery or by
exercise. Each open contract has one long and one short; a trade raises
open interest when both sides open, lowers it when both close, and leaves it
unchanged when one side hands its position to the other.
\end{definition}

\begin{definition}[Price limit]\label{def:m1:futures-contracts-and-exchanges:limit}
On a future, a \emph{\omterm{def:m1:asian-and-emerging-equity-markets:limit}{price limit}} (for shares,
\cref{def:m1:asian-and-emerging-equity-markets:limit}) is a bound, set by the
exchange relative to a reference price, beyond which the contract may not
trade during a session or a part of it; reaching it may trigger a halt.
\end{definition}

\begin{dated}{2026-09}{How ES settles each day and when it stops}\label{dat:m1:futures-contracts-and-exchanges:limits}
The daily settlement of the lead-month E-mini is the volume-weighted average
price of trades between 14:59:30 and 15:00:00 Chicago time, with the
midpoint of the bid and ask as fallback; other months are settled from
calendar-spread prices. Downward \omterm{def:m1:asian-and-emerging-equity-markets:limit}{price limits} of 7\%, 13\% and 20\% from a
reference price apply during US stock-market hours, coordinated with the
stock market's circuit breakers (\cref{ch:m1:exchange-traded-funds}), and a
symmetric band applies overnight (5\% either way in the 2016 rule filing
consulted; read the current rule).
\end{dated}

Since all months of a product settle at three o'clock and a position is
marked there, the last thirty seconds before the settlement are, for a
futures desk, what the closing auction is for an equity desk.

Volume moves from one expiry to the next in a few days before expiry: the
\emph{roll}. Holders who want to keep their exposure sell the \omterm{def:m1:futures-contracts-and-exchanges:code}{front month} and
buy the next as a calendar spread; \omterm{def:m1:futures-contracts-and-exchanges:oi}{open interest} migrates, and the \omterm{def:m1:futures-contracts-and-exchanges:code}{front
month} of the data changes. A price series stitched across rolls without
adjustment contains jumps that are not returns
(\cref{ch:m1:basis-roll-and-delivery}).

\begin{omfigure}
\begin{tikzpicture}
\begin{axis}[width=11cm,height=4.8cm,
  xlabel={trading days to expiry of the front contract},ylabel={open interest (million)},
  tick label style={font=\small},label style={font=\small},
  xmin=-30,xmax=9,ymin=0,ymax=2.8,grid=major,grid style={gray!25},
  legend columns=2,legend style={font=\footnotesize,at={(0.5,-0.32)},anchor=north,draw=none,fill=none,/tikz/every even column/.append style={column sep=8pt}},
  table/col sep=comma]
\addplot[omDef,very thick] table[x=day,y=front_m]{figdata/markets-1/18-futures-contracts-and-exchanges/roll.csv};
\addplot[omThm,very thick,dashed] table[x=day,y=next_m]{figdata/markets-1/18-futures-contracts-and-exchanges/roll.csv};
\legend{front month,next month}
\end{axis}
\end{tikzpicture}

\omcaption{A stylised quarterly roll: \omterm{def:m1:futures-contracts-and-exchanges:oi}{open interest} passes from the front
contract to the next within about a week, centred here eight trading days
before expiry. The date and speed differ by product and must be measured,
not assumed. Data: the tutorial's simulation.}\label{fig:m1:futures-contracts-and-exchanges:roll}
\end{omfigure}

\section{Who trades futures}

\begin{definition}[Hedger and speculator]\label{def:m1:futures-contracts-and-exchanges:hedger}
A \emph{hedger}\index{hedger} holds futures to offset a risk it has
elsewhere: a producer's crop, a fund's shares, a dealer's bonds. A
\emph{speculator}\index{speculator} holds them for the exposure itself. The
US regulator's weekly report on positions classifies large traders as
\emph{commercial}, those who use the contract for hedging as its regulation
defines it, and \emph{non-commercial}.
\end{definition}

The textbook pair hides the population that matters to a trading firm:
\omterm{def:m1:what-a-trading-firm-does:asset-manager}{asset managers} equitising cash and managing duration; leveraged funds in
basis, trend and relative-value trades; dealers hedging swaps and options;
and \omterm{def:m1:what-a-trading-firm-does:market-maker}{market makers}, who hold little overnight and account for much of the
volume. A future is also the hedge of choice for every \omterm{def:m1:delta-one-instruments:deltaone}{delta-one} instrument
of \cref{ch:m1:delta-one-instruments}, which ties its price to the cash market
within the bands of \cref{ch:m1:basis-roll-and-delivery}.

\begin{method}[Sizing a futures hedge]\label{met:m1:futures-contracts-and-exchanges:hedge}
To hedge a portfolio of value $V$ and beta $\beta$ to the index with a future
at price $F$ and multiplier $m$: sell $N = \beta V/(mF)$ contracts, rounded
to the nearest integer. The residual exposure is $\beta V - NmF$, at most half
a contract; for small portfolios use the smaller contract if one exists.
Beta is an estimate: its error is usually far larger than the rounding.
\end{method}

\begin{omfigure}
\begin{tikzpicture}
\begin{axis}[width=11cm,height=4.8cm,
  xlabel={portfolio value (\$ million)},ylabel={unhedged residual (\%)},
  tick label style={font=\small},label style={font=\small},
  xmin=0.1,xmax=3,ymin=0,ymax=55,grid=major,grid style={gray!25},
  legend columns=2,legend style={font=\footnotesize,at={(0.5,-0.32)},anchor=north,draw=none,fill=none,/tikz/every even column/.append style={column sep=8pt}},
  table/col sep=comma]
\addplot[omDef,thick] table[x=portfolio_m,y=residual_pct_es]{figdata/markets-1/18-futures-contracts-and-exchanges/residual.csv};
\addplot[omThm,thick,dashed] table[x=portfolio_m,y=residual_pct_micro]{figdata/markets-1/18-futures-contracts-and-exchanges/residual.csv};
\legend{contract of \$300\,000,contract one tenth the size}
\end{axis}
\end{tikzpicture}

\omcaption{Rounding to whole contracts: the residual exposure of a
beta-one portfolio, in percent of its value. Below \$150\,000 the large
contract cannot hedge at all (residual 100\%, off the scale). Data:
\cref{met:m1:futures-contracts-and-exchanges:hedge}.}\label{fig:m1:futures-contracts-and-exchanges:residual}
\end{omfigure}

\section{Tutorial: a contract table}\label{tut:m1:futures-contracts-and-exchanges:table}

\begin{tutorial}
\textbf{Goal.} Hold six contract specifications in one structure and derive
notional, \omterm{def:m1:futures-contracts-and-exchanges:tick}{tick value}, \omterm{def:m1:futures-contracts-and-exchanges:tick}{tick} in basis points, leverage and hedge sizes from
it. \textbf{End state:} the three data figures of this chapter.
\begin{tutsteps}
\item \textbf{One record per contract.} Bond multipliers are expressed per
price point: \$100\,000 of face value quoted in percent is \$1\,000 a point.
\omcode{code/markets-1/18-futures-contracts-and-exchanges/python/futures_basics.py}{8}{27}{A specification and the three numbers derived from it.}\label{lst:m1:futures-contracts-and-exchanges:spec}
\item \textbf{Hedge sizing.} Round to the nearest contract and report what
is left.
\omcode{code/markets-1/18-futures-contracts-and-exchanges/python/futures_basics.py}{44}{48}{Contracts to sell and the residual exposure.}\label{lst:m1:futures-contracts-and-exchanges:hedge}
\item \textbf{Leverage.} Divide notional by the \omterm{def:m1:clearing-and-settlement:margin}{initial margin} your broker
quotes today (\cref{ch:m1:margin}); with a margin of 5\% of notional a 2\% move
is 40\% of the deposit.
\end{tutsteps}
\textbf{What to change next.} Add four contracts from the exchanges'
product pages, a short-term interest-rate future among them, whose price is
100 minus a rate and whose \omterm{def:m1:futures-contracts-and-exchanges:tick}{tick value} comes from a day-count, not from a
price.
\end{tutorial}

\section{Build: the contract master}\label{bld:m1:futures-contracts-and-exchanges:master}

\begin{build}
\bfield{Purpose} Every component of the miniature firm that touches a
future (the P\&L keeper of \cref{ch:m1:pnl-and-positions}, the margin engine
of \cref{ch:m1:margin}, order entry) asks one module what the contract is.
\bfield{Interface} \texttt{ContractMaster.from\_csv(path)};
\texttt{spec(root)} with exact \texttt{multiplier}, \texttt{tick\_size},
\texttt{tick\_value}, \texttt{to\_ticks(price)}, \texttt{from\_ticks(n)},
\texttt{notional(price)}; \texttt{parse(symbol, today)} returning an
instrument; \texttt{expiry(instrument)};
\texttt{parse\_thirty\_seconds(text)}.
\bfield{Rules} All arithmetic in exact fractions; a price off the \omterm{def:m1:futures-contracts-and-exchanges:tick}{tick} grid
is an error, never rounded. A single-digit year resolves to the first such
year that is not in the past. A month code not listed for the product is an
error. Expiry rules that need an exchange calendar raise rather than guess.
\bfield{Data} \texttt{data/markets-1/contracts\_sample.csv}, six rows
transcribed from this chapter's sources: a test fixture, not reference data.
\bfield{Acceptance tests} \texttt{code/firm/contracts/tests/}: exact \omterm{def:m1:futures-contracts-and-exchanges:tick}{tick
values}; round trips through integer \omterm{def:m1:futures-contracts-and-exchanges:tick}{ticks}; Treasury quotes; symbols and
third Fridays; validation.
\bfield{Stretch} Load an exchange's daily reference-data file and diff it
against yesterday's: tick-size and listing changes are events
(\cref{ch:m1:reading-market-data}).
\end{build}

\begin{omsources}
\item CME, Rulebook Chapter 358 (E-mini S\&P 500 futures) and Chapter 351,
as filed with the CFTC (March 2016); CME Group, daily settlement procedures
for equity-index futures, as filed with the CFTC (January 2018).
\item CBOT, Rulebook Chapter 19 (Treasury note futures), and the CFTC filing
of 21 February 2025 on deliverable grades; NYMEX, Rulebook Chapter 200.
\item Eurex, product pages \emph{EURO STOXX 50 Index Futures} and
\emph{Euro-Bund Futures}; ICE Futures Europe, \emph{Brent Crude Futures}
product page.
\item US Commodity Futures Trading Commission, \emph{Commitments of Traders:
Explanatory Notes}.
\end{omsources}

\section{Exercises}

\begin{exercise}[$\star$]\label{exo:m1:futures-contracts-and-exchanges:1}
A trader is long 12 E-minis from 6\,012.50; the price falls to 5\,987.25. Give
the move in \omterm{def:m1:futures-contracts-and-exchanges:tick}{ticks} and the loss.
\end{exercise}

\begin{exercise}[$\star$]\label{exo:m1:futures-contracts-and-exchanges:2}
The 10-year note future goes from \texttt{111-245} to \texttt{112-080}. Give
both prices in decimal, the move in \omterm{def:m1:futures-contracts-and-exchanges:tick}{ticks} and the profit on one long
contract.
\end{exercise}

\begin{exercise}[$\star$]\label{exo:m1:futures-contracts-and-exchanges:3}
Today is 18 September 2026. Decode \texttt{FESXH7} and give its last trading
day. What would \texttt{ESH6} mean today?
\end{exercise}

\begin{exercise}[$\star\star$]\label{exo:m1:futures-contracts-and-exchanges:4}
From \cref{fig:m1:futures-contracts-and-exchanges:tickbp}, the \omterm{def:m1:futures-contracts-and-exchanges:tick}{tick} of FESX is
more than four times that of ES relative to value. Give two consequences
for a \omterm{def:m1:what-a-trading-firm-does:market-maker}{market maker} and one for a fund crossing the spread on \euro50 million.
\end{exercise}

\begin{exercise}[$\star\star$]\label{exo:m1:futures-contracts-and-exchanges:5}
A broker asks \$21\,000 of \omterm{def:m1:clearing-and-settlement:margin}{initial margin} per E-mini at 6\,000. Give the
leverage and the index move that loses the whole deposit. Why is the
deposit not the right measure of the money needed to hold the position?
\end{exercise}

\begin{exercise}[$\star\star$]\label{exo:m1:futures-contracts-and-exchanges:6}
A new contract lists. Trades: A buys 5 from B; C buys 3 from A; B buys 2
from C; D buys 4 from A. Give the \omterm{def:m1:futures-contracts-and-exchanges:oi}{open interest} and each trader's position
after each trade. What were volume and \omterm{def:m1:futures-contracts-and-exchanges:oi}{open interest} at the end?
\end{exercise}

\begin{exercise}[$\star\star\star$]\label{exo:m1:futures-contracts-and-exchanges:7}
\emph{Coding.} From \texttt{residual.csv}: for portfolios of \$1 million and
more, report the largest residual with the \$300\,000 contract and with the
contract one tenth its size. State the general bound and check it.
\end{exercise}

\begin{exercise}[$\star\star\star$]\label{exo:m1:futures-contracts-and-exchanges:8}
\emph{Find the flaw.} A commentary reads: ``\omterm{def:m1:futures-contracts-and-exchanges:oi}{Open interest} in the front
contract rose 20\% today while the price rose: new buyers are pouring in, and
buyers now outnumber sellers.'' Correct it, and say what the same data can
legitimately tell you.
\end{exercise}

\section{Problem: Sizing a Hedge}

\begin{problem}[{Weekend problem --- a transatlantic portfolio and two futures}]\label{pb:m1:futures-contracts-and-exchanges:1}
On 18 September 2026 a fund wants to remove the market exposure of two
sleeves for three months: \$187 million of US shares with a beta of 1.15 to
the S\&P 500, and \euro43 million of euro-area shares with a beta of 0.90 to
the EURO STOXX 50. The E-mini is at 6\,000 and the FESX at 5\,400.

\medskip\noindent\textbf{Part I --- The contracts.}
\begin{enumerate}
\item Give the notional of one contract of each.
\item Give the exact and the rounded number of E-minis to sell, and the
residual.
\item Same for the FESX.
\item Which expiry should the fund sell, with its symbol and last trading
day, if it does not want to roll?
\item With \omterm{def:m1:clearing-and-settlement:margin}{initial margin} at 5\% of notional, give the deposit for the US
hedge.
\end{enumerate}

\medskip\noindent\textbf{Part II --- A bad day.}
The S\&P 500 falls 2\%.
\begin{enumerate}[resume]
\item Give the gain on the futures and the expected loss on the US sleeve.
\item When does the cash from the futures arrive, and when is the loss on
the shares realised? Why does the difference matter on a \emph{good} day?
\item The sleeve actually loses 2.9\%. Give the hedged P\&L and name its
source.
\item The fund's beta estimate has a standard error of 0.10. Compare the
corresponding exposure with the rounding residual.
\end{enumerate}

\medskip\noindent\textbf{Part III --- Details.}
\begin{enumerate}[resume]
\item In December the fund extends the hedge by a quarter and rolls to
March. The calendar spread has a \omterm{def:m1:futures-contracts-and-exchanges:tick}{tick} of 0.05. Give the cost of paying one
spread \omterm{def:m1:futures-contracts-and-exchanges:tick}{tick} on the whole position.
\item The European hedge is in euros and the fund reports in dollars. What
exposure has the hedge \emph{not} removed?
\item The E-mini's daily settlement is a thirty-second average at 15:00
Chicago time, which is 16:00 in New York; the shares are valued at their
closing auction at the same moment. What appears in the daily P\&L?
\item Why might the fund hedge the US sleeve with ES even if most of it is
small-capitalisation stocks, and what does it give up?
\item The hedge will be in place for three months. Using
\cref{ch:m1:delta-one-instruments}, what does it cost or earn besides
commissions?
\end{enumerate}

\medskip\noindent\textbf{Part IV --- Judgement.}
\begin{enumerate}[resume]
\item After the hedge, what risks does the fund still run?
\item A 7\% fall triggers the first \omterm{def:m1:asian-and-emerging-equity-markets:limit}{price limit} during US hours. What happens
to the hedge, and to the shares?
\item Would a contract one tenth the size improve the hedge here?
\item Why do funds hedge with futures and not by selling the shares?
\item State the \emph{named result}: the number of E-minis and the residual
exposure.
\item In one sentence: what does a \omterm{def:m1:futures-contracts-and-exchanges:future}{futures contract} standardise, and what
does it leave to the trader?
\end{enumerate}
\end{problem}

\section{Interview questions}

\begin{interviewq}[$\star$ \iqroles{trader, developer, researcher}]\label{iq:m1:futures-contracts-and-exchanges:1}
What is the \omterm{def:m1:futures-contracts-and-exchanges:tick}{tick value} of an E-mini S\&P 500 future, and what is the
contract worth?
\end{interviewq}

\begin{interviewq}[$\star$ \iqroles{trader, researcher}]\label{iq:m1:futures-contracts-and-exchanges:2}
What is \omterm{def:m1:futures-contracts-and-exchanges:oi}{open interest}, and how does it differ from volume?
\end{interviewq}

\begin{interviewq}[$\star\star$ \iqroles{developer, researcher}]\label{iq:m1:futures-contracts-and-exchanges:3}
How would you represent futures prices in a trading system, and why?
\end{interviewq}

\begin{interviewq}[$\star\star$ \iqroles{trader, researcher}]\label{iq:m1:futures-contracts-and-exchanges:4}
Why does a future trade on only one exchange when a share trades on many
venues?
\end{interviewq}

\begin{interviewq}[$\star\star$ \iqroles{researcher, mle}]\label{iq:m1:futures-contracts-and-exchanges:5}
You are given ten years of daily front-month prices. What must you do before
computing returns?
\end{interviewq}

\begin{interviewq}[$\star\star\star$ \iqroles{trader, researcher}]\label{iq:m1:futures-contracts-and-exchanges:6}
You must hedge \$200 million of equities for a quarter. Walk me through every
decision.
\end{interviewq}
